\documentclass[a4paper, 12pt]{article}

%%% SST LAB PROTOCOLL PREAMBLE
%%% 2019
%%%%%%%%%%%%%%%%%%%%%%%%%%%%%%%


%%% PACKAGES
%%%%%%%%%%%%%%%%%%%%%%%%%%%

\usepackage[ngerman]{babel}

\usepackage[utf8]{inputenc}
\usepackage{amsmath}
\usepackage{pgfplots}
\usepackage{tikz}
\usepackage[many]{tcolorbox}
\usepackage{graphicx}
\graphicspath{ {./graphics/} }
\usepackage{pdfpages}
\usepackage{dashrule}
\usepackage{float}
\usepackage{siunitx}
\usepackage{trfsigns}
\usepackage{booktabs}
\usepackage[european]{circuitikz}
\usepackage{tcolorbox}

%%% DOCUMENT GEOMETRY
%%%%%%%%%%%%%%%%%%%%%%%%%%%

\usepackage{geometry}
\geometry{
 a4paper,
 total={0.6180339887498948\paperwidth,0.6180339887498948\paperheight},
 top = 0.1458980337503154\paperheight,
 bottom = 0.1458980337503154\paperheight
 }
\setlength{\jot}{0.013155617496424828\paperheight}
\linespread{1.1458980337503154}

\setlength{\parskip}{0.013155617496424828\paperheight} % paragraph spacing


%%% COLORS
%%%%%%%%%%%%%%%%%%%%%%%%%%%

\definecolor{red1}{HTML}{f38181}
\definecolor{yellow1}{HTML}{fce38a}
\definecolor{green1}{HTML}{95e1d3}
\definecolor{blue1}{HTML}{66bfbf}
\definecolor{hsblue}{HTML}{00b1db}
\definecolor{hsgrey}{HTML}{afafaf}

%%% CONSTANTS
%%%%%%%%%%%%%%%%%%%%%%%%%%%
\newlength{\smallvert}
\setlength{\smallvert}{0.0131556\paperheight}


%%% COMMANDS
%%%%%%%%%%%%%%%%%%%%%%%%%%%

% differential d
\newcommand*\dif{\mathop{}\!\mathrm{d}}

% horizontal line
\newcommand{\holine}[1]{
  	\begin{center}
	  	\noindent{\color{hsgrey}\hdashrule[0ex]{#1}{1pt}{3mm}}\\%[0.0131556\paperheight]
  	\end{center}
}

% mini section
\newcommand{\minisec}[1]{ \noindent\underline{\textit {#1} } \\}

% quick function plot
\newcommand{\plotfun}[3]{
  \vspace{0.021286\paperheight}
  \begin{center}
    \begin{tikzpicture}
      \begin{axis}[
        axis x line=center,
        axis y line=center,
        ]
        \addplot[draw=red1][domain=#2:#3]{#1};
      \end{axis}
    \end{tikzpicture}
  \end{center}
}

% box for notes
\newcommand{\notebox}[1]{

\tcbset{colback=white,colframe=green1!100!black,title=Note!,width=0.618\paperwidth,arc=0pt}

 \begin{center}
  \begin{tcolorbox}[]
   #1 
  \end{tcolorbox}
 
 \end{center} 
 
}

% box for equation
\newcommand{\eqbox}[2]{
	
	\tcbset{colback=white,colframe=green1!100!black,title=,width=#2,arc=0pt}
	
	\begin{center}
		\begin{tcolorbox}[ams align*]
				#1
		\end{tcolorbox}
		
	\end{center} 
	
}
% END OF PREAMBLE

%%%%%%%%%%%%%%%%%%%%%%%%%%%%%%%%%%%%%

\begin{document}

%%%%%%%%%%%%%%%%%%%%%%%%%%%%%%%%%%%%%
  \includepdf{./titlepage/titlepage.pdf}
  \clearpage
  \setcounter{page}{1}
%%%%%%%%%%%%%%%%%%%%%%%%%%%%%%%%%%%%%
\section{Tasterprellen}
Das Prellen der Taster ist mechanischen Ursprungs. Ursache ist das elastische
Zurückprallen gegen die Federung (Wikipedia), wodurch der Kontakt oft mehrmals
geöffnet und geschlossen wird bevor der eigentlich stabile Tasterzustand
entsteht, was dazu führen kann, dass fehlerhaft mehrfache Anschläge
registriert werden.

Hardwaremäßig kann man ein Tiefpassfilter einsetzen um die kurzen,
ungewollten Impulse zu filtern. Nachfolgend wird dann ein Schmitt-Trigger zur
Erzeugung der korrekten Signalwerte (High/Low) geschaltet.

In der Software kann man eine bestimmte Entprellzeit abwarten, bevor der
Tasterzustand gelesen wird. Dort wird z.B. der Interrupt durch das erste Prellen
ausgelöst, der Zustand registriert und mit dem Zustand nach der Entprellzeit verglichen.

Tasterprellzeiten liegen im Bereich 100 $\, \si{\micro\second}$ bis 10 $\, \si{\milli\second}$.

\section{Codierung des Tastendrucks}
Mithilfe einer integrierten Encoderschaltung, wie dem 74HC148, kann ein Tastendruck
aus acht Tastern in einen 3-bit Code gewandelt werden. Jeder Taster erhält also
einen eigenen Code, z.B. 101, über den er von den anderen unterscheidbar ist.
Werden mehrere Taster gleichzeitig gedrückt, liegt u.U. eine Priosierungslogik
vor, die den Taster/Code mit der höheren Priorität einstellt und den/die mit der
niedrigeren damit unterdrückt.

Der Group Select (GS) Pin am 74HC148 wird bei Eingangsänderung gesetzt und kann
zur Interrupterzeugung verwendet werden.

\section{Interruptkonfiguration}

\subsubsection*{Global Interrupt Enable (GIE) im SR muss aktiviert sein}

\begin{lstlisting}
__bis_SR_register(GIE);
\end{lstlisting}

\subsubsection*{ Port Interrupt Enable (PxIE) muss aktiviert sein}
\begin{lstlisting}
P1IE |= BIT4;
\end{lstlisting}

\subsubsection*{Interrupt Edge Select (IES) wählen}
\begin{lstlisting}
P1IES |=  BIT4; // Interrupt bei High-Low-Flanke
P1IES &=~ BIT4; // Interrupt bei Low-High-Flanke
\end{lstlisting}

\subsubsection*{Interrupt Service Routine, Vektor zuweisen}
\begin{lstlisting}
#pragma vector=PORT1_VECTOR
__interrupt void port1_isr_name() {
    // ISR
}
\end{lstlisting}

\section{Tastendruck einlesen}
\begin{itemize}
\item Portinitialisierung, Interruptkonfiguration
\item ISR: Ausgelöst durch GS
\item 3-bit Wert aufnehmen
\item Entprellzeit abwarten
\item Entsprechend Zustand setzen, ob Taster gedrückt
\item Weiterverarbeitung außerhalb der ISR abhängig vom Tasterzustand
\end{itemize}



  
 \end{document}

% Local Variables:
% TeX-engine: luatex
% End:
